\section{Legal Foundations}

This appendix summarizes the key Supreme Court decisions that informed the Invocation of Rights. Each case reflects a specific vulnerability in civilian–police encounters—where silence, hesitation, or ambiguous language has been held insufficient to preserve constitutional protections. The cases listed below support the legal design of the script’s four lines, as detailed in Section 4.

\begin{longtable}{@{}P{0.22\textwidth} P{0.46\textwidth} P{0.30\textwidth}@{}}
\toprule
\textbf{Case (Year)} & \textbf{Holding} & \textbf{Script Relevance}\\
\midrule
\endfirsthead
\toprule
\textbf{Case (Year)} & \textbf{Holding} & \textbf{Script Relevance}\\
\midrule
\endhead
\midrule
\multicolumn{3}{r}{\emph{Continued on next page}}\\
\midrule
\endfoot
\bottomrule
\endlastfoot

\textit{Berghuis v.\ Thompkins (2010)} &
Remaining silent is not enough to invoke the Fifth Amendment; it must be clearly stated. &
Justifies: “I invoke my right to remain silent.”\\[4pt]

\textit{Salinas v.\ Texas (2013)} &
Pre‑custodial silence can be used against a suspect unless the Fifth Amendment is explicitly invoked. &
Reinforces the need for clear invocation even before arrest.\\[4pt]

\textit{Davis v.\ United States (1994)} &
An ambiguous reference to counsel (“Maybe I should…”) does not require police to stop questioning. &
Justifies: “I invoke my right to a lawyer.”\\[4pt]

\textit{Edwards v.\ Arizona (1981)} &
Once a suspect clearly requests a lawyer, police must stop questioning unless the suspect re‑initiates. &
Reinforces legal force of an explicit lawyer request.\\[4pt]

\textit{Schneckloth v.\ Bustamonte (1973)} &
Police are not required to inform civilians they can refuse a search; unclear responses may count as consent. &
Supports: “I don’t consent to any searches.”\\[4pt]

\textit{Florida v.\ Bostick (1991)} &
Passive agreement or compliance can be treated as valid consent during police questioning. &
Reinforces the need for clear verbal refusal.\\[4pt]

\textit{Florida v.\ Royer (1983)} &
Officers turned a consensual encounter into an unlawful detention without cause. &
Justifies: “I want to leave. Am I free to go?”\\[4pt]

\textit{Terry v.\ Ohio (1968)} &
Police may conduct brief investigative stops with reasonable suspicion, but not without limits. &
Clarifies the boundary this line seeks to test.\\[4pt]

\textit{Hiibel v.\ Sixth Judicial District (2004)} &
States may require individuals to provide identification during lawful stops. &
Notes the Invocation’s compatibility with stop‑and‑identify statutes.\\
\end{longtable}

\begin{raggedright}
\emph{Note:} The cases listed reflect federal constitutional standards. While state law may expand protections, the rights asserted in the script—grounded in Supreme Court precedent—are recognized nationwide. The Invocation is designed to function regardless of custody status or jurisdiction, offering civilians a consistent legal foundation across encounters.
\end{raggedright}
